% =============================================================================
% CVE-2026-60004 — Xpectra Security Research
% Assembled by build.sh: cat main.tex part2.tex part3.tex part4.tex > paper-full.tex
% Every number here comes from designs/FACTS.md or a cited evidence file.
% =============================================================================
% !TEX program = pdflatex
\documentclass[journal]{IEEEtran}
\usepackage[utf8]{inputenc}
\usepackage{graphicx}
\usepackage{booktabs}
\usepackage{array}
\usepackage{xcolor}
\usepackage{listings}
\usepackage{hyperref}
\usepackage{microtype}
\usepackage{subcaption}

\hypersetup{colorlinks=true,linkcolor=blue,citecolor=blue,urlcolor=blue}

\definecolor{codebg}{RGB}{245,247,250}
\definecolor{codekw}{RGB}{0,90,150}
\definecolor{codecm}{RGB}{120,120,120}
\definecolor{codestr}{RGB}{160,40,40}

\lstdefinestyle{gosrc}{
  language=Go,
  backgroundcolor=\color{codebg},
  basicstyle=\ttfamily\scriptsize,
  keywordstyle=\color{codekw}\bfseries,
  commentstyle=\color{codecm}\itshape,
  stringstyle=\color{codestr},
  frame=single,
  breaklines=true,
  breakatwhitespace=false,
  columns=fullflexible,
  keepspaces=true,
  showstringspaces=false,
  tabsize=2,
  xleftmargin=0.5em, xrightmargin=0.3em
}
\lstdefinestyle{diffsrc}{
  language={},
  backgroundcolor=\color{codebg},
  basicstyle=\ttfamily\scriptsize,
  commentstyle=\color{codecm}\itshape,
  frame=single,
  breaklines=true,
  breakatwhitespace=false,
  columns=fullflexible,
  keepspaces=true,
  showstringspaces=false,
  tabsize=2,
  xleftmargin=0.5em, xrightmargin=0.3em
}
\lstdefinestyle{transcript}{
  language={},
  backgroundcolor=\color{codebg},
  basicstyle=\ttfamily\scriptsize,
  frame=single,
  breaklines=true,
  breakatwhitespace=false,
  columns=fullflexible,
  keepspaces=true,
  showstringspaces=false,
  xleftmargin=0.5em, xrightmargin=0.3em
}

\newcommand{\figdir}{figures}

\begin{document}

\title{Beyond the Published Exploit: Process-Level Proof, Live-Negative Fix
Verification, and the Detection Consequences of Gitea's \texttt{diffpatch} RCE
(CVE-2026-60004)}

% JOURNAL MODE: \IEEEauthorblockN/A are inline dummies in journal mode — stack manually.
\author{Xpectra Security Research\\[0.5ex]
\small Autonomous Red/Blue Operations}

\markboth{Beyond the Published Exploit: CVE-2026-60004}{Xpectra Security
Research \MakeLowercase\expandafter\the\year}

\maketitle

\begin{abstract}
Gitea's \texttt{diffpatch} API endpoint allowed arbitrary shell execution as
the service operating-system account (CVE-2026-60004, CWE-94, affected
\texttt{>=1.17}, \texttt{<1.27.1}). The vendor advisory, published 2026-07-28,
did not merely describe the flaw: it shipped a working Python exploit of
roughly 330 lines, including an exfiltration technique that needs no outbound
channel. \textbf{This paper makes no discovery claim.} It contributes what that
advisory left unproven: process-level proof of execution captured from an
independent reimplementation (parent process \texttt{git write-tree}, hook
process \texttt{/bin/sh\hspace{0pt}hooks/post-index-change\hspace{0pt}0\hspace{0pt}0},
working directory inside the bare temporary clone, \texttt{GIT\_DIR=.},
\texttt{uid=1000(git)}); empirical verification of the three preconditions the
vendor only asserts, including the finding that all three hold in the official
Docker image without setting any precondition-related option (the only
container parameters were standard hostname, port, and install-lock values
equivalent to a standard install's own, none part of the precondition set);
an A/B fix verification whose
negative arm is live rather than vacuous (22 of 22 cycles without execution,
with arrival of the payload at the sink path positively observed in 2 of 22
sampled cycles, and the second \texttt{diffpatch} response changing from HTTP 201
to HTTP 500); a measured
impact radius (\texttt{app.ini} of 1323 bytes with 3 secret-bearing lines read,
planted secrets read back, arbitrary write under \texttt{/data/gitea}
confirmed, and one negative result); and an analysis of the fix whose decisive
hunk changes one boolean argument and leaves the \texttt{git apply} flags
untouched --- so a byte-identical patch is inert after the fix, which makes
patch-content-based detection structurally weak and forces behavioural
detection. All experiments ran in an isolated, zero-egress Docker lab against a
fictional internal domain; no third-party system was touched. CISA added the
CVE to the KEV catalog on 2026-08-25 with a federal deadline of 2026-08-28.
\end{abstract}

\begin{IEEEkeywords}
Gitea, CVE-2026-60004, remote code execution, Git hooks, vulnerability
verification, detection engineering, reproducible experiments.
\end{IEEEkeywords}

% =============================================================================
\section{Introduction}
\label{sec:intro}

A source-hosting service that applies attacker-supplied patches has to decide
where those patches land before they are ever committed. Gitea 1.27.0 applied
them inside a \emph{bare} temporary clone, and in a bare repository the
checkout root \emph{is} \texttt{\$GIT\_DIR}. A patch that adds an executable
file at \texttt{hooks/post-index-change} therefore does not drop a stray file:
it installs a live Git hook that Git then executes while writing the index.
That is CVE-2026-60004 --- arbitrary command execution as the Gitea service
user, from one authenticated, write-capable API call, sent twice.

This CVE is unusual in a way that shapes every sentence that follows:
\textbf{the vendor published the working exploit inside the advisory}
(GHSA-rcr6-4jqh-j84m, published 2026-07-28T17:34:46Z). The advisory carries the
full mechanism, the vulnerable command line, a complete Python
proof-of-concept of roughly 330 lines, and an exfiltration technique that
stores command output in Git objects and retrieves it as a branch through
authenticated smart HTTP --- no outbound connection required. We make no
discovery claim anywhere in this paper, in our video material, or in any
companion artifact; Section~\ref{sec:record} states the division of credit
explicitly.

What an advisory cannot do for you is the rest of the security work:
\emph{showing that it actually happens, showing that the fix actually stops
it, and showing what a defender can actually see}. Working in an isolated
zero-egress Docker lab (\texttt{gitea/gitea:1.27.0} against \texttt{1.27.1},
one variable at a time) with a canary-first payload that writes one file and
does nothing else, we contribute:

\begin{enumerate}
\item \textbf{Process-level proof of execution}, which the advisory does not
provide. Our hook reports its own lineage: parent process
\texttt{/usr/bin/git write-tree}, own command line
\texttt{/bin/sh\hspace{0pt}hooks/post-index-change\hspace{0pt}0\hspace{0pt}0}, working
directory
\texttt{/data/gitea/tmp/\allowbreak local-repo/\allowbreak upload.git749828539},
\texttt{GIT\_DIR=.}, \texttt{uid=1000(git)}, plus the container's process tree
captured from inside the hook (Section~\ref{sec:repro}). The public record says
generically that Git invokes the hook ``while writing the index''; in the public
sources we checked (advisory, CVE record, fix commit message, patch diff) the
parent process is not named.
\item \textbf{Empirical verification of preconditions the vendor only asserts},
and a derived finding that changes triage: in the official image as published,
all three hold \emph{without setting any precondition-related option}
(Section~\ref{sec:precond}).
\item \textbf{Fix verification with a live negative}, not an absence of
evidence. On \texttt{1.27.1} with byte-identical payload: no execution in 22 of
22 cycles, the second \texttt{diffpatch} submission changed from HTTP 201 to
HTTP 500, and a polling monitor still observed the payload materialising at the
sink path --- so the arm proves the payload arrives and still does not run
(Section~\ref{sec:fix}).
\item \textbf{The non-obvious consequence of the fix}, which is the most useful
thing here for a defender. The decisive hunk changes
\texttt{Clone(\dots,\allowbreak true)} to \texttt{Clone(\dots,\allowbreak false)}
and adds no path deny-list, no mode filter, and no change to the
\texttt{git apply} flags. The same patch bytes are still accepted and still land
on disk after the fix. Consequence: patch-content detection cannot distinguish
an attack that executed from one that was neutralised, and an artifact rule
fires on both (Sections~\ref{sec:fixchange} and~\ref{sec:detect}).
\item \textbf{Measured impact radius} in place of an enumerated list:
\texttt{app.ini} readable (1323 bytes, 3 secret-bearing lines), planted secrets
read back, arbitrary write confirmed where the service lives, repositories
visible --- and one honest negative, the parent process environment leaked
nothing (Section~\ref{sec:impact}).
\item \textbf{A negative-results log}. Four of our own measurements were wrong
or unprovable and they are printed in Section~\ref{sec:negative}, including a
\texttt{grep} that appeared to prove the \texttt{diffpatch} route was absent
because it inspected a 581-byte wrapper instead of the 124\,MB binary. A
published exploit does not make the surrounding evidence trustworthy by
itself.
\end{enumerate}

Section~\ref{sec:mech} gives the mechanism plus an independent static
corroboration; Section~\ref{sec:repro} the reproduction and process evidence;
Section~\ref{sec:fix} the A/B; Section~\ref{sec:detect} what a defender can
actually see; Section~\ref{sec:remediation} remediation and incident response;
Section~\ref{sec:limitations} limitations and negative results; and
Section~\ref{sec:ethics} ethics, scope and reproducibility. Every number was
either measured in our lab or read from the public record with a date. Nothing
here is estimated.

% =============================================================================
\section{The Public Record, Ethics Position, and Timeline}
\label{sec:record}

\subsection{What the vendor published}

Table~\ref{tab:record} separates the two bodies of work. Nothing in this paper
moves from the right column to the left.

\begin{table*}[t]
\centering
\footnotesize
\begin{tabular}{@{}p{7.7cm}p{8.1cm}@{}}
\toprule
\textbf{Public record (vendor and registries)} & \textbf{This work} \\
\midrule
Mechanism: add/add collision, three-way fallback, bare root as
\texttt{\$GIT\_DIR}, live hook (GHSA ``Details''). &
Process-level proof of that mechanism: parent PID and command line, hook
command line, \texttt{cwd}, \texttt{GIT\_DIR}, resolved hook path, container
process tree, wall-clock timing. \\
Vulnerable command \texttt{git apply --index --recount --cached --binary} plus
\texttt{-3} on Git $\geq$2.32, in
\texttt{services/repository/files/patch.go}. &
Independent static corroboration of the collision detector recovered from the
shipped binary (\texttt{git ls-files -u -z}, Section~\ref{sec:mech}), and the
measurement that Git is 2.54.0 in both images. \\
Working exploit \texttt{gitea\_diffpatch\_rce\_poc.py} ($\approx$330 lines) with
output-free exfiltration via Git objects plus a branch fetched over
authenticated smart HTTP. &
A deliberately weaker stdlib-only canary PoC: no outbound channel, no reverse
shell, no destructive action; the hook writes one canary file with process
evidence (Section~\ref{sec:repro}). We did \emph{not} reproduce the vendor's
exfiltration channel. \\
Three stated preconditions: Git $\geq$2.32, enabled \texttt{diffpatch} route,
writable and executable temporary filesystem. &
Those three measured rather than assumed, plus the derived finding that the
official image satisfies all three by default
(Section~\ref{sec:precond}). \\
Impact stated as a list: application secrets, database credentials, mounted
repositories, OAuth and integration credentials. &
Impact measured: byte counts, secret-bearing line counts, planted secrets read
back, arbitrary-write confirmation, network visibility, and one negative
(Section~\ref{sec:impact}). \\
Fix commit \texttt{d7bc52be\ldots} (2026-07-26) and first patched version
1.27.1. &
A/B execution of the fix with a live negative, plus a response-code change
(HTTP 500 on the second submission) that the advisory does not mention
(Section~\ref{sec:fix}). \\
\bottomrule
\end{tabular}
\caption{Credit division. Everything in the left column belongs to the vendor
and the registries; we cite it and do not claim it. The right column is what
this paper adds.}
\label{tab:record}
\end{table*}

\subsection{Timeline and operational urgency}
\label{sec:timeline}

\begin{figure*}[t]
\centering
\includegraphics[width=\textwidth]{\figdir/fig-timeline.png}
\caption{Disclosure and response chronology, read from the registries rather
than remembered: fix commit merged 2026-07-26T17:32:02Z (three days
\emph{before} the advisory), GHSA published 2026-07-28T17:34:46Z and updated
18:40:54Z the same day, CISA KEV \texttt{dateAdded} 2026-08-25 with FCEB
\texttt{dueDate} 2026-08-28, and NVD publication 2026-08-26T20:17:56.010Z
with status \emph{Analyzed}. The KEV note references BOD 26-04, whose
staggered windows are 3, 14 and 60 days.}
\label{fig:timeline}
\end{figure*}

Two facts in Figure~\ref{fig:timeline} make this urgent rather than academic.
First, the fix reached the repository on 2026-07-26, three days before the
advisory: an operator who tracks commits had a patch before a disclosure
notice. Second, the response clock arrived a month later and is short --- CISA
added the CVE to the KEV catalog on 2026-08-25, 28 days after the advisory,
NVD published it the following day with \emph{Analyzed} status, and the federal
(FCEB) deadline is 2026-08-28, three days after the KEV entry. In this
case the operational urgency preceded the NVD record by a day. The KEV entry's own notes reference BOD 26-04, with staggered
windows of 3, 14 and 60 days.

Two catalog measurements put that deadline in proportion. The KEV record
shipped with this paper was fetched on 2026-08-27 (catalog
\texttt{2026.08.27}, 1685 entries in 1\,618\,253 bytes, full-file SHA-256
pinned in the supporting material): 1672 entries (99.2\,\%) already had a
passed \texttt{dueDate} and only 13 were due within 30 days. A KEV entry is a
prioritised signal, not a self-enforcing control. The same entry records
\texttt{knownRansomwareCampaignUse = Unknown}, so we do not claim ransomware
use, and NVD records no CVSS v4.0 score, so we do not quote one. The severity
of record is the NVD-measured CVSS 3.1 vector
\texttt{AV:N/\hspace{0pt}AC:L/\hspace{0pt}PR:N/\hspace{0pt}UI:N/\hspace{0pt}S:U/\hspace{0pt}C:H/\hspace{0pt}I:H/\hspace{0pt}A:H}
= 9.8 (Critical), typed \emph{Secondary}, with CWE-94.

\subsection{Scope statement}

Every experiment ran in a Docker lab we created and destroyed: an
\texttt{--internal} bridge with verified zero egress, a fictional hostname, no
real users and no real data. Payloads were canary-first by policy.
Section~\ref{sec:ethics} gives the full ethics, consent and reproducibility
statement rather than repeating it here.
% =============================================================================
\section{Mechanism, and an Independent Static Corroboration of It}
\label{sec:mech}

\subsection{The chain as executed}

\begin{figure*}[t]
\centering
\includegraphics[width=\textwidth]{\figdir/fig-chain.png}
\caption{The executed chain. An authenticated \texttt{POST diffpatch} is
applied with \texttt{git apply --index --recount --cached --binary -3} into a
shared \emph{bare} temporary clone under \texttt{/data/gitea/tmp/local-repo/}.
Submitting the \emph{same} patch twice produces an add/add collision; Git's
three-way fallback checks the path out into the working tree even though the
operation was requested with \texttt{--cached}. In a bare repository the
checkout root \emph{is} \texttt{\$GIT\_DIR}, so \texttt{hooks/post-index-change}
with mode \texttt{100755} is a live hook, and Git executes it while writing the
index. The dashed box is what the vendor publishes and we deliberately did
\emph{not} reproduce: output exfiltration through Git objects plus a branch
fetched over authenticated smart HTTP. Our PoC stops at the hook writing one
canary file.}
\label{fig:chain}
\end{figure*}

Four properties of this chain are individually reasonable and jointly fatal.

\textbf{(1) The apply command is not the bug.} The command line published in
the advisory is what Gitea 1.27.0 runs (Listing~\ref{lst:apply}).
\texttt{--cached} means ``update the index, not the working tree''; on a normal
repository that is an ordinary, safe instruction.

\begin{lstlisting}[style=gosrc,
caption={The patch-application command of Gitea 1.27.0, as published in the
\texttt{Details} section of the advisory
(\texttt{services/repository/files/patch.go}).},
label={lst:apply}]
cmdApply := gitcmd.NewCommand("apply", "--index", "--recount",
    "--cached", "--binary")
if git.DefaultFeatures().CheckVersionAtLeast("2.32") {
    cmdApply.AddArguments("-3")
}
\end{lstlisting}

\textbf{(2) The collision is the surprise.} Applying the same
\texttt{new file} patch twice turns the second application into an add/add
conflict, and the three-way path (\texttt{-3}) resolves it by checking the
indexed path out into the working tree --- contradicting \texttt{--cached} in
exactly the one case where the difference is exploitable. This is not a parser
bug or a memory error; it is two documented behaviours composed.

\textbf{(3) Bare makes the root equal \texttt{\$GIT\_DIR}.} In the temporary
clone, checkout root and Git directory are the same directory, so the
checked-out path \texttt{hooks/post-index-change} is not a stray project file:
it is \texttt{\$GIT\_DIR/hooks/post-index-change}, a hook Git is designed to
execute. Our measurement confirms this operationally rather than by code
reading: the hook reported \texttt{GIT\_DIR=.} with the temporary clone as its
working directory (Listing~\ref{lst:canary}).

\textbf{(4) The response lies by omission.} The hook's exit status is not
propagated to the \texttt{diffpatch} response. Both of our submissions returned
HTTP 201 with a commit SHA while the payload executed side-effectfully, as the
advisory also states. The operational consequence: a defender watching API
responses sees two successful, unremarkable commits. Response-code monitoring is
not a detection strategy here; only process or artifact telemetry sees the
attack.

\subsection{Independent static corroboration from the shipped binary}

The advisory's mechanism claim has a static footprint that we did not see
mentioned in any public source. The \texttt{1.27.0} binary carries the format
string of Listing~\ref{lst:conflict}. \texttt{git ls-files -u} lists
\emph{unmerged} index entries --- the collision detector of a three-way merge.
Its presence in this handler confirms, without running anything, that the code
path expects unmerged entries, which is what the add/add mechanism produces.
Corroboration is not proof of execution; Section~\ref{sec:repro} is the proof.

\begin{lstlisting}[style=transcript,
caption={Format string recovered from \texttt{/app/gitea/gitea} in the
\texttt{gitea/gitea:1.27.0} image (our measurement; not named in the advisory).
\texttt{git ls-files -u} is the unmerged-entry query of a three-way merge.},
label={lst:conflict}]
conflict detected at: git ls-files -u -z: %w
\end{lstlisting}

\subsection{The smallest credential the attack needs}
\label{sec:precond-cred}

The advisory states that open registration is required only for the
no-prior-credentials path. We measured that path instead of assuming it.
\texttt{GET /user/sign\_up} returned HTTP 200. A \texttt{POST /user/sign\_up}
carrying only the four fields that the real form exposes (\texttt{user\_name},
\texttt{email}, \texttt{password}, \texttt{retype}), with an empty
\texttt{\_csrf} and no prior authenticated cookie, returned HTTP 303 and created
an account that the instance's own administrative CLI listed as
\texttt{IsActive true}, \texttt{IsAdmin false}. Every experiment in this paper
uses that kind of account: non-administrative, with write access only to a
repository it created itself. No administrator, no token, no complicit user,
zero clicks.

% =============================================================================
\section{Preconditions the Vendor Asserts --- Measured}
\label{sec:precond}

The advisory lists three trigger requirements and no verification data. We
measured each against the images as published (Table~\ref{tab:precond}).

\begin{table}[t]
\centering
\footnotesize
\begin{tabular}{@{}p{2.2cm}p{5.4cm}@{}}
\toprule
Precondition (vendor wording) & Measurement and verdict \\
\midrule
Git $\geq$2.32 & \texttt{git version 2.54.0} inside \texttt{gitea/gitea:1.27.0}
\emph{and} \texttt{1.27.1}. Satisfied. \\
An enabled \texttt{diffpatch} route & The real binary
\texttt{/app/gitea/gitea} (124\,MB) contains
\texttt{/repos/\%s/\%s/diffpatch}, \texttt{ApplyRepoDiffPatch},
\texttt{routers/\allowbreak api/\allowbreak v1/\allowbreak repo.\allowbreak ApplyDiffPatch}. Satisfied. \\
Writable \emph{and} executable temp filesystem & A script written to
\texttt{/tmp} inside the container executed successfully
(\texttt{TMP\_EXEC\_OK}); \texttt{/tmp} is \texttt{rw} overlay without
\texttt{noexec}. Satisfied. \\
\bottomrule
\end{tabular}
\caption{The three vendor preconditions, measured. All satisfied without
modifying the deployment.}
\label{tab:precond}
\end{table}

\textbf{Derived finding, and the triage argument it removes.} These are not
properties of a hardened or unusual deployment. They are properties of the
official Docker image as the vendor publishes it, with only the standard
environment variables for hostname, port and installation lock, and no other
configuration touched. All three hold simultaneously out of the box, which
removes the standard deflection --- ``this only fires if you configured it
badly''. For triage, the operative question is not \emph{whether} an instance
meets the preconditions but \emph{whether} it runs below 1.27.1, and
\texttt{GET /api/v1/version} answers that without credentials.

% =============================================================================
\section{Reproduction with Process-Level Proof}
\label{sec:repro}

\subsection{Lab}

Two containers on one \texttt{--internal} Docker bridge
(\texttt{cve60004\_lab}, \texttt{172.31.224.0/24}) with no external route:
\texttt{gitea-vuln} (\texttt{gitea/gitea:1.27.0},
\texttt{172.31.224.2:3000}) and \texttt{gitea-patched} (\texttt{1.27.1},
\texttt{172.31.224.3:3000}), each with \texttt{/data} on a named volume and a
fictional hostname. The attacker is the host at the bridge gateway
(\texttt{172.31.224.1}); verification of results is \emph{out of band}
(\texttt{docker exec} on the target), because our PoC has no output channel by
design. Zero egress was verified, not assumed: from inside the container,
\texttt{wget} to \texttt{1.1.1.1} reports \emph{Network unreachable} and to
\texttt{api.github.com} reports a dead resolver (\emph{bad address}).
Table~\ref{tab:timing} comes from \texttt{exploit/cve-2026-60004-canary.py}, a
stdlib-only reimplementation that shares no code with the vendor's PoC.
Fig.~\ref{fig:run} is the operator's screen during the filmed run that produced the
evidence in this section.

\begin{figure*}[t]
\centering
\includegraphics[width=\textwidth]{\figdir/fig-still-run-landing.png}
\caption{One recorded take of the operator's screen at the instant the reverse
shell lands: two terminals stacked, captured simultaneously. The upper pane is
the external listener: \texttt{CHANNEL OPEN} carrying the run nonce and
\texttt{stage=rshell}, the landing identity \texttt{uid=1000(git)}, and the
listener's own accept-to-channel timing line; the agent's first two commands
answer \texttt{whoami $\rightarrow$ git} and \texttt{id $\rightarrow$
uid=1000(git)}. The lower pane is the vulnerable service driving it: the state
machine at \texttt{STATE 5 REVERSE\_SHELL} through the double \texttt{diffpatch}
(\texttt{scratch repo} \texttt{labattacker/f4-s5-291012 $\rightarrow$ HTTP 201}).
The \texttt{getcwd} errors in the upper pane are real: the hook ran from a
directory the double apply had deleted, and cutting them would imply a run we
did not have. The still is taken from the recorded take that the published cut
uses with zero speed changes; paper stills carry no burned subtitles --- those
are the film's typography, not the screen's content. No
on-screen count is quoted from it --- the strings shown are in
\texttt{evidence/chain-show/} and the alert counts in
\texttt{evidence/rule-fires/}. Frame provenance and crop geometry are in
\texttt{paper/FIGURES.md}.}
\label{fig:run}
\end{figure*}

\subsection{The payload we chose, and why}

The vendor's PoC executes an arbitrary command and exfiltrates its output
through Git objects. \textbf{We did not reproduce that channel.} Our hook writes
one canary file inside the target and nothing else: no network access, no
shell, no change to service state. That costs nothing evidentially, because the
claim under test is ``attacker-controlled content executes as the service
user'', and a file containing a nonce planted before the attack, written by a
process the attacker did not otherwise control, proves precisely that. The
canary also produced the process evidence the advisory does not contain
(Listing~\ref{lst:canary}).

\begin{lstlisting}[style=transcript,
caption={Process-level proof, read out of band from the target after the attack
(\texttt{evidence/canary-content.txt}, a 940-byte file; long lines are wrapped
here for column width). Top block: the hook reporting its own lineage --- the
executed nonce, \texttt{uid=1000(git)}, its own PID 361 and command line, and
parent 360. Bottom block: the container process tree captured \emph{by the hook
itself}, showing \texttt{s6-supervise gitea} $\rightarrow$ \texttt{gitea web}
$\rightarrow$ \texttt{git write-tree} $\rightarrow$ the hook shell. The
truncated \texttt{executed\_at\_utc} fraction is a busybox \texttt{date}
limitation documented in Section~\ref{sec:negative}; the proof is the nonce and
the lineage, not the timestamp.},
label={lst:canary}]
CANARY_EXEC_OK cve-2026-60004
nonce=c643d0b3-2d58-4177-837e-ab19db8b681f
executed_at_utc=2026-08-27T14:33:22.
whoami=git uid=1000 gid=1000
self_pid=361
parent_pid=360
parent_cmdline=/usr/bin/git write-tree
self_cmdline=/bin/sh hooks/post-index-change 0 0
cwd=/data/gitea/tmp/local-repo/
  upload.git749828539
git_dir_env=.
hook_resolved=/data/gitea/tmp/local-repo/
  upload.git749828539/hooks/post-index-change

--- proc tree view ---
PID   PPID  USER     COMMAND
    1     0 root     /usr/bin/s6-svscan /etc/s6
   13     1 root     s6-supervise openssh
   14     1 root     s6-supervise gitea
   15    13 root     sshd: /usr/sbin/sshd -D -e [listener] 0 of 10-100
   16    14 git      /usr/local/bin/gitea web
  347    16 git      /usr/bin/git cat-file --batch-command
  360    16 git      /usr/bin/git write-tree
  361   360 git      {post-index-chan} /bin/sh
                      hooks/post-index-change 0 0
  371   361 git      ps -o pid,ppid,user,args
\end{lstlisting}

\textbf{Three items the public record does not contain.} (i) The parent process
is \texttt{/usr/bin/git write-tree}. The advisory says only that Git invokes the
hook ``while writing the index''; naming the parent matters because
``\texttt{git write-tree} with a shell child'' is an anomaly a defender can write
a rule against, whereas ``a hook ran'' is not. (ii) \texttt{git\_dir\_env=.}
turns the architectural claim ``the bare root is \texttt{\$GIT\_DIR}'' into an
observed environment variable. (iii) The hook resolves inside a temporary upload
repository whose name is \texttt{upload.git} followed by digits --- the artifact
signature the attack leaves behind, and the path a defender should watch.

\subsection{Wall-clock cost of the attack}
\label{sec:timing}

\begin{table}[t]
\centering
\footnotesize
\begin{tabular}{@{}p{2.0cm}p{1.0cm}p{1.0cm}p{1.0cm}p{1.7cm}@{}}
\toprule
Step & Run A & Run B & Run C & Response \\
\midrule
Auth (non-admin) & $0.02$\,s & $0.03$\,s & $0.02$\,s & HTTP 200 \\
Create repository & $0.30$\,s & $0.32$\,s & $0.29$\,s & HTTP 201 \\
\texttt{diffpatch} 1/2 & $0.57$\,s & $0.60$\,s & $0.55$\,s & HTTP 201 \\
\texttt{diffpatch} 2/2 & $0.86$\,s & $0.89$\,s & $0.82$\,s & HTTP 201 \\
\midrule
Total delivery & $\mathbf{0.86}$\,s & $\mathbf{0.89}$\,s & $\mathbf{0.82}$\,s & codes $[201,201]$ \\
\bottomrule
\end{tabular}
\caption{Delivery timing measured by the PoC over three complete runs on
\texttt{1.27.0}. The columns are bound to their artifacts explicitly: run A is
\texttt{evidence/poc-run-2.txt} (0.86\,s total), run B is
\texttt{evidence/poc-run-1.txt} (0.89\,s), run C is
\texttt{evidence/poc-run-current-english.txt} (0.82\,s), and the published range is
the minimum and maximum of those three totals. The step rows are \emph{cumulative from
the start of the run}, not per-step durations: they do not add up to the total, which equals
the last step because delivery completes with it. Runs A and B were recorded with an
earlier build of the script whose console
labels were in Spanish; run C was recorded with the shipped English build, which
also adds the scope-guard line. The payload bytes are unaffected by either
change. Both \texttt{diffpatch} responses report
success; neither reveals that a hook executed. The script's own output states
that no result of the script is proof by itself, and the canary was verified
separately, out of band.}
\label{tab:timing}
\end{table}

The entire delivery --- self-registered non-admin authentication, private
repository creation, two \texttt{diffpatch} submissions carrying the same bytes
--- costs $0.82$--$0.89$\,s from the attacker position across three runs, against
two API routes that a source-hosting service must expose to be useful. Combined
with property~(4) of Section~\ref{sec:mech}, the attack leaves no anomalous
API-level trace. We measured that trace: the service's own router log records
the two \texttt{201} responses for the attack and 41 internal
\texttt{pre-receive}/\texttt{post-receive} callbacks proving a push happened,
and \textbf{zero} lines mention \texttt{post-index-change}, the hook that
actually executed; with the stock \texttt{[log]} configuration no request log is
written to \texttt{/data/gitea/log} at all
(\texttt{evidence/app-log-visibility.txt}). The log view of this exploit is two
successful commits.

% =============================================================================
\section{Measured Impact Radius, Including a Negative}
\label{sec:impact}

The advisory lists what exploitation \emph{may} expose. We probed through the
same primitive (a second canary-style payload, \texttt{exploit/impact-probe.py},
delivered over the identical two-request chain) and report what it \emph{did}
expose, with sizes and counts. Every secret read was either planted by us for
this purpose or is the service's own configuration; the lab contains no
third-party data to read.

\begin{table}[t]
\centering
\footnotesize
\begin{tabular}{@{}p{2.2cm}p{5.4cm}@{}}
\toprule
Probe & Measured result \\
\midrule
Service config & \texttt{app.ini}
(\texttt{\allowbreak /data/\allowbreak gitea/\allowbreak conf/app.ini}) readable: \textbf{1323 bytes}; the
probe counted \textbf{3} secret-bearing lines, and they are \texttt{SECRET\_KEY}
(present but \emph{empty} in this lab), \texttt{INTERNAL\_TOKEN} and
\texttt{JWT\_SECRET}. The \texttt{DATABASE} section is readable too, but its
\texttt{PASSWD} is empty here, so it is listed as exposure surface, not as a
credential we obtained \\
Planted secrets & \texttt{DB\_PASSWORD} and \texttt{AWS\_ACCESS\_KEY\_ID}
(fictional lab values) read back in clear \\
Repository tree & client repositories listable under
\texttt{/data/git/repositories/} \\
Arbitrary write & \texttt{ESCRITURA\_OK} on
\texttt{\allowbreak /data/\allowbreak gitea/\allowbreak attacker\_write\_\allowbreak test.txt}, deleted after the probe \\
Process environment & \textbf{No} variable containing
\texttt{SECRET}, \texttt{TOKEN} or \texttt{PASS} in the parent process
environment \\
Network from the service context & \texttt{lo} and \texttt{172.31.224.2/24}
only (zero-egress lab) \\
\bottomrule
\end{tabular}
\caption{Impact probe results, read out of band
(\texttt{evidence/impact-probe-content.txt}). Rows are what the payload
demonstrated, not what it could theoretically reach. Token strings such as
\texttt{ESCRITURA\_OK} are quoted \emph{verbatim} from the captured output: the shipped
probe prints its console labels in Spanish, which \texttt{evidence/PROVENANCE.md} records
explicitly, and translating a token here would separate the table from its artifact.}
\label{tab:impact}
\end{table}

\textbf{Reading the table honestly.} The confidentiality row is the expensive
one in a real deployment: a readable \texttt{app.ini} containing
\texttt{INTERNAL\_TOKEN} and \texttt{JWT\_SECRET} hands the attacker the signing and
token-verification material of the source host, which turns a source-hosting
compromise into a code-integrity and credential problem for everything that trusts
those tokens. Database credentials are a further \emph{inference}, not a measurement:
this lab backs Gitea with SQLite and leaves \texttt{PASSWD} empty, so no database
credential was read in this experiment. Where an \texttt{app.ini} names a real
PostgreSQL or MySQL password, the same file readability is the ordinary path to those
credentials, which is why remediation treats them as exposed. The write row is
what makes this more than a read bug: arbitrary write as \texttt{uid=1000(git)}
inside the volume that holds configuration and repositories. The environment row
is a genuine negative and we keep it in the paper --- the advisory's impact list
includes process-environment secrets, and in this deployment shape that line
\emph{did not reproduce}.

\textbf{A forensic constraint with detection consequences.} The executed hook
lived in a temporary upload repository below
\texttt{/data/gitea/tmp/local-repo/}, which is short-lived. The artifact that
proves compromise therefore \emph{disappears on its own}: post-incident disk
forensics is not a reliable path to this attack, and detection has to be
behavioural and contemporaneous (Section~\ref{sec:detect}).
% =============================================================================
\section{Fix Verification with a Live Negative}
\label{sec:fix}

\subsection{What the fix actually changes}
\label{sec:fixchange}

The fix commit is public: \texttt{d7bc52be\allowbreak eadff4be5f5690de4d5de42a\allowbreak bd10affe},
dated 2026-07-26T17:32:02Z, titled ``refactor: git patch apply (\#38637)
(\#38638)'', touching \texttt{cherry\_pick.go}, \texttt{patch.go},
\texttt{patch\_test.go} and \texttt{update.go}. Inside that refactor sits one
decisive hunk (Listing~\ref{lst:fix}). The regression test added in the same
commit asserts only that \texttt{.git} exists inside the temporary repository:
the test encodes ``the patch sandbox must not be bare'' and nothing more.

\begin{lstlisting}[style=diffsrc,
caption={The decisive hunk of the fix (\texttt{services/repository/files/patch.go};
full diff in \texttt{designs/fix-commit-full.diff}). One boolean turns the patch
sandbox from bare to non-bare. The \texttt{git apply} flags of
Listing~\ref{lst:apply} are not touched.},
label={lst:fix}]
-  if err := t.Clone(ctx, opts.OldBranch, true); err != nil {
+  // here must NOT use bare repo, because the following git
+  // commands might operate working tree ("--index") directly
+  if err := t.Clone(ctx, opts.OldBranch, false); err != nil {
       return nil, err
  }
\end{lstlisting}

\textbf{The consequence that is easy to miss.} The fix changes \emph{where} a
patch is applied. It does not change the \texttt{git apply} flags; it adds no
path deny-list for \texttt{hooks/}, no filter on executable modes, and does not
disable the three-way fallback. The collision mechanism is still there --- the
sandbox simply no longer makes its checkout root double as
\texttt{\$GIT\_DIR}. Two operational consequences follow.

\begin{enumerate}
\item \textbf{A byte-identical exploit is inert, and still looks the same.} The
exploit bytes are not rejected, not sanitised, and not made anomalous; they are
applied to a worktree, where \texttt{hooks/post-index-change} is a plain file.
We confirmed this by observation rather than inference: on the patched instance
the monitor still saw the hook file materialise on disk
(Table~\ref{tab:ab}).
\item \textbf{Content-signature detection is therefore structurally weak.} A
signature over the patch text --- a \texttt{hooks/} path, \texttt{new file mode
100755}, the \texttt{diff} header --- fires identically on a vulnerable and a
patched instance, and fires on any legitimate patch that touches hook files. It
can raise an attempt signal; it cannot tell an operator whether they were
compromised, and it cannot tell them whether they are patched. What \emph{can}
discriminate is behaviour: process lineage, and which directory Git considered
\texttt{\$GIT\_DIR}. That is the most actionable claim in this paper, and it
comes from reading the fix rather than the exploit.
\end{enumerate}

\subsection{A/B experiment}

Same lab network, same configuration, same account shape, byte-identical payload;
the only variable is the image (\texttt{1.27.0} against \texttt{1.27.1}). The
negative arm is not accepted on absence of evidence: a polling monitor
(\texttt{exploit/ab-monitor.sh}) searched for the materialised hook file under
\texttt{/data/gitea/tmp/local-repo/} throughout each patched run, so that ``no
execution'' cannot be quietly explained by ``the payload never arrived''.

\begin{table}[t]
\centering
\footnotesize
\begin{tabular}{@{}p{2.7cm}p{1.9cm}p{2.9cm}@{}}
\toprule
Observation & \texttt{1.27.0} vuln. & \texttt{1.27.1} patched \\
\midrule
\texttt{diffpatch} 1/2 & HTTP 201 & HTTP 201 \\
\texttt{diffpatch} 2/2 & HTTP 201 & \textbf{HTTP 500} \\
Hook file at the sink path & yes & \textbf{observed} on disk, in a temporary
\texttt{upload.git*} repository \\
Canary created & yes & \textbf{no --- 22/22 cycles} \\
Hook content executed & yes, \texttt{uid=1000} & none observed \\
Delivery window & $0.82$--$0.89$\,s & not applicable \\
\bottomrule
\end{tabular}
\caption{A/B result. Identical payload and configuration. The patched arm is a
\emph{live} negative: the payload reaches the sink path and is still not
executed, and the second submission now fails with HTTP 500 instead of returning
a successful commit.}
\label{tab:ab}
\end{table}

\begin{figure}[t]
\centering
\includegraphics[width=\columnwidth]{\figdir/fig-still-vulnerable-arm.png}
\caption{The victim repository on the \emph{vulnerable} arm, mid-defacement on
the filmed run: the page's own footer reads \texttt{Version:~1.27.0} and the
address bar reads \texttt{172.31.224.2}; the crop keeps the full footer band,
its geometry measured in \texttt{paper/FIGURES.md}. The README has been
overwritten through the same primitive that landed the shell --- the table
visible here is the agent's own impact note, declaring which value was
replaced and how the original bytes are restored. The patched arm is not shown
on camera --- the demonstration film was ruled offensive-only --- so this
figure identifies only the arm it shows, and the patched-arm case rests on the
A/B table above and \texttt{evidence/fix-verification/}, not on a frame.}
\label{fig:arm}
\end{figure}

\textbf{Why the negative is non-vacuous.} Three observations have to hold
together for the patched arm to mean anything, and all three do. Fig.~\ref{fig:arm}
shows the arm under test as the operator saw it. (i) The payload
reaches the target and is written: the monitor captured
\texttt{hooks/post-index-change} under a temporary upload repository on the
\emph{patched} instance. (ii) The failure is visible in the protocol: the second
submission returns HTTP 500 where the vulnerable instance returns HTTP 201.
(iii) Nothing executes across 22 independent cycles --- we ran the chain 22 times
rather than once, because a single negative proves only that one run did not
fire. Given a non-bare sandbox, the checked-out hook lands in
\texttt{<worktree>/hooks/}, which is not \texttt{\$GIT\_DIR/hooks/}, so Git has no
reason to execute it. That is exactly the failure mode the hunk in
Listing~\ref{lst:fix} predicts, and that is the sense in which the fix is
verified: not ``it fails'', but ``it fails in the way, and only in the way, that
the change implies''.

We do not claim to have identified which internal call produces the HTTP 500; we
report the observed status code, not a mechanism we did not measure.

\subsection{Severity, checked against what we could prove}
\label{sec:cvss}

The severity of record is the NVD CVSS 3.1 assessment
\texttt{AV:N/AC:L/PR:N/UI:N/S:U/C:H/I:H/A:H} = 9.8 (Critical), typed
\emph{Secondary}, CWE-94. Our measurements are consistent with it: the endpoint
is network-reachable, no privilege beyond a self-created account is needed while
registration is open (Section~\ref{sec:precond-cred}), no user interaction is
involved, and total delivery took under a second --- so \texttt{AC:L} is not a
generous reading. Two honesty notes. First, \texttt{PR:N} rides entirely on open
registration: with registration disabled the prerequisite becomes one ordinary
authenticated account, effectively \texttt{PR:L}, and the vulnerability remains
exploitable --- the bar is ``any user who can create a repository'', not ``an
administrator''. Second, we did not test availability impact; our payload policy
forbids destructive action, so \texttt{A:H} is an inference from arbitrary write
access as the service user, not a result we produced.

% =============================================================================
\section{What a Defender Can Actually See}
\label{sec:detect}

This section states detection \emph{requirements} derived from measured
artifacts, and it quantifies the shipped rule package (\texttt{rules/}) against
captures we recorded ourselves. We publish no detection rate and no false-positive
rate that we did not measure: every figure below is re-derivable from the fire
reports in \texttt{evidence/rule-fires/}. What we can state is which signals
discriminate and which do not --- and that distinction is the useful part.
The sensor itself is evidenced by the fire reports and the tables below, not by
a screenshot: what we publish per signal is the command that ran, the artifact
it consumed, and the alert lines it produced.

\subsection{Measured behaviour of the shipped signatures}

\textbf{What counts.} The package carries seven CVE SIDs, of which \emph{six}
alert: \texttt{sid:202660045} is a \texttt{flowbits:\,set} carrying
\texttt{noalert}, so it can never appear in an alert count, and a harness that
reports ``seven signatures fired'' is wrong by construction. Replayed through
Suricata~8.0.6, the attack capture yields \textbf{11 alerts} from the six
alerting SIDs (\texttt{41:2}, \texttt{42:2}, \texttt{43:2}, \texttt{44:1},
\texttt{46:2}, \texttt{47:2}); both captures are \texttt{suricata -r} replays of
recorded wire data, not live sensor output.

\textbf{False positives, measured.} A benign negative is only worth what its
traffic is worth. Our simple capture (\texttt{wire-benign-control.pcap},
48\,packets) contains no request with the shape of the attack, so its zero
CVE-SID alerts prove that the engine decoded the bytes and nothing about
specificity --- reading it as a specificity result is the error we corrected
while preparing this paper. We therefore recorded a structurally similar control,
\texttt{wire-benign-structural.pcap} in that same folder (123 packets read),
whose traffic is what a legitimate administrator generates: three
\texttt{diffpatch} cycles that add an executable \texttt{maintenance/helper.sh},
built with the same patch generator the shipped PoC uses, plus ordinary
\texttt{git-upload-pack}\allowbreak/\allowbreak\texttt{git-receive-pack} traffic.

On that benign traffic the mode-bit signature \texttt{sid:202660041} fired on
\textbf{3 of 3} legitimate cycles, while the five signatures keyed to the
\texttt{hooks/} target path produced no alert. One of the six alerting
rules, on one capture, is not a measurement of false-positive rate in the field;
it is a measurement of how little a mode bit discriminates. That result is why
\texttt{sid:202660041} ships at \texttt{priority:3} with
\texttt{confidence~low} and its measured rate recorded in rule metadata, and why
we do not recommend it as a standalone alert.

\subsection{Signals that do not discriminate}

\textbf{The patch content.} Two \texttt{POST} requests to the
\texttt{diffpatch} route carrying identical bodies, the second declaring a
\texttt{new file mode 100755} under \texttt{hooks/}, are the shape of this
attack on the wire, and our capture (\texttt{evidence/wire-attack.pcap}, 33\,400
bytes, 72 packets, filtered on \texttt{tcp port 3000}) contains it alongside
benign control traffic --- a \texttt{GET /} and one ordinary \texttt{POST
contents} --- recorded precisely so that false-positive behaviour can be
measured. But per Section~\ref{sec:fixchange} the same bytes are accepted and
written on a patched instance, so a content signature is an \emph{attempt}
indicator with no bearing on outcome.

\textbf{The HTTP response.} Both \texttt{diffpatch} calls return HTTP 201 on a
vulnerable instance while code executes, because the hook's exit status is not
propagated. Success-code monitoring sees nothing. Note the asymmetry we
measured: it is the \emph{patched} instance that answers 500, so a ``many 5xx on
\texttt{diffpatch}'' alert is a statement about the fix, not about compromise.

\textbf{Disk forensics after the fact.} The hook lives in a short-lived
temporary upload repository, so the artifact that proves execution is gone by
the time an operator looks (Section~\ref{sec:impact}).

\subsection{Signals that do discriminate}

\begin{enumerate}
\item \textbf{Process lineage.} The parent of the hook shell is
\texttt{/usr/bin/git write-tree} and the child is \texttt{/bin/sh
hooks/post-index-change 0 0}, under the \texttt{gitea web} process, as
\texttt{uid=1000} (Listing~\ref{lst:canary}). A \texttt{git} command that writes
a tree should not have a shell descendant. This is the highest-fidelity signal we
measured: across the 394 service-log lines covering our run, not one mentions the
executing hook while 40 name the \texttt{diffpatch} endpoint itself, so the
request is visible and the execution is not
(\texttt{evidence/app-log-visibility.txt}); and it fires only when execution
actually happened.
\item \textbf{Location, not existence, of the hook file.} A file named
\texttt{hooks/post-index-change} inside a directory that \emph{is}
\texttt{\$GIT\_DIR} is compromise; the same file inside a worktree is inert. The
discriminating observable is therefore \texttt{GIT\_DIR=.} together with the
temporary path (\texttt{/data/gitea/tmp/local-repo/}, name prefix
\texttt{upload.git}), not the filename. A rule keyed on the filename alone
produces the same alert on patched and vulnerable hosts.
\item \textbf{Repeated identical \texttt{diffpatch} against one repository.} The
exploit \emph{requires} the same patch twice inside one second
(Table~\ref{tab:timing}). That repetition is a behavioural invariant of the
attack, independent of the file it carries.
\item \textbf{A new ref in a repository with no push.} This is the vendor's
retrieval channel: output committed as an object, surfaced as a branch, fetched
over authenticated smart HTTP. We did not reproduce it, so we list it as a
requirement the advisory implies rather than a signal we validated.
\end{enumerate}

\subsection{MITRE ATT\&CK mapping of what we demonstrated}

Table~\ref{tab:attack} is deliberately narrow: every row names a technique and the
concrete observation that backs it in this paper, with a section reference. Rows we
could not back are absent rather than inferred. In particular there is no Lateral
Movement row and no Exfiltration row, because we did not reproduce the vendor's
output-free retrieval channel and did not move laterally in the lab; we write a
\emph{requirement} the advisory implies in Section~\ref{sec:detect} instead of a
technique we observed. Mapping is at tactic level only: we make no claim about which
sub-technique an intrusion would be attributed to in the field.

\begin{table}[t]
\centering
\footnotesize
\begin{tabular}{@{}p{1.3cm}p{2.0cm}p{4.2cm}@{}}
\toprule
Tactic & Technique & Evidence in this paper \\
\midrule
Initial Access & T1190 & Self-registered non-admin account; HTTP 303 on
\texttt{sign\_up} (\S\ref{sec:precond-cred}) \\
Execution & T1059.004 & \texttt{/bin/sh hooks/}\allowbreak\texttt{post-index-change
0 0}, \texttt{uid=1000(git)} (Listing~\ref{lst:canary}) \\
Discovery & T1083 & \texttt{app.ini} and the repository tree enumerated from
inside the hook (\S\ref{sec:impact}) \\
Credential Access & T1552.001 & 3 secret-bearing lines in \texttt{app.ini};
planted secrets read back \\
Persistence & T1546.004 & Git hook installed through repository-controlled
content (lab only) \\
\bottomrule
\end{tabular}
\caption{Techniques mapped only to behaviour we actually observed. Lateral
Movement and Exfiltration are absent on purpose: the vendor's output-free
exfiltration channel was not reproduced here.}
\label{tab:attack}
\end{table}
% =============================================================================
\section{Remediation and Incident Response}
\label{sec:remediation}

\textbf{The fix is the remediation.} Upgrade to Gitea 1.27.1 or later: the
affected range is \texttt{>=1.17, <1.27.1} and 1.27.1 is the first patched
release. Version is verifiable unauthenticated through
\texttt{GET /api/v1/version}, and our precondition measurements
(Section~\ref{sec:precond}) mean that an operator should not spend time
checking whether their configuration is ``bad enough'' --- in the reference
Docker deployment all three vendor preconditions hold by default. Because
CISA added this CVE to the KEV catalog on 2026-08-25 with a federal deadline
of 2026-08-28, and because the advisory itself ships a working exploit, an
internet-reachable Gitea below 1.27.1 should be treated as already targeted,
not hypothetically vulnerable.

\textbf{What an interim mitigation can and cannot do.} Removing open
registration raises the attacker's prerequisite from \emph{none} to
\emph{one ordinary account} --- it does not mitigate the flaw, because the
vulnerability triggers from ordinary write access to any repository
(Section~\ref{sec:precond-cred}). Do not treat registration policy, WAF rules
on patch content, or response-code alerting as substitutes: Sections
\ref{sec:fixchange} and \ref{sec:detect} explain why each of those is weak
against this specific fix shape. If an upgrade must wait, restrict who can
create repositories and who can reach \texttt{diffpatch}, and deploy the
process-lineage detection of Section~\ref{sec:detect} so that an attempt
becomes visible at the only layer that discriminates.

\textbf{Incident response, keyed to measured artifacts.}
\begin{enumerate}
\item \emph{Triage (read-only).} Enumerate running processes for a
\texttt{git} command with a shell descendant, specifically
\texttt{git write-tree} $\rightarrow$ \texttt{/bin/sh\allowbreak
hooks/}\ldots; list files under
\texttt{/data/gitea/tmp/\allowbreak local-repo/\allowbreak upload.git*/hooks/}
(the window is short, so do this continuously rather than on demand); review
reverse-proxy or access logs for repeated \texttt{diffpatch} POSTs to the same
repository inside a one-second window, and for 500s on that route.
\item \emph{Assume credential exposure proportional to the measured impact.}
In our lab the hook read a 1323-byte \texttt{app.ini} whose three secret-bearing
lines are \texttt{SECRET\_KEY} (empty here), \texttt{INTERNAL\_TOKEN} and
\texttt{JWT\_SECRET}, and wrote an arbitrary file under \texttt{/data/gitea}. No
database credential was obtained in this experiment (the lab uses SQLite with an
empty \texttt{PASSWD}), but a readable \texttt{app.ini} is the ordinary path to one
wherever the database has a real password, so if execution is confirmed, rotate
application secrets and database credentials,
and treat repository contents and any integration tokens reachable from that
configuration as exposed --- not because we demonstrated misuse, but because
we demonstrated readability and writability.
\item \emph{Eradication and recovery.} Rebuild rather than clean a host where
arbitrary write by the service user is confirmed inside the configuration
volume. Upgrade first, then verify the fix \emph{as a behaviour}: our A/B
protocol (Section~\ref{sec:fix}) is 22 cycles of an inert canary payload with
an arrival monitor --- the same three lines of evidence (no canary, HTTP 500
on the second submission, payload still observed on disk) are what an operator
should require before declaring an instance fixed. Re-running the vendor PoC
is not a fix test: it proves only that the exploit ``did not appear to work''
on one run.
\end{enumerate}

% =============================================================================
\section{Limitations and Negative Results}
\label{sec:limitations}

This section is the part of the paper an advisory-shaped summary would drop.
It stays in, because four of our own measurements were wrong or unprovable,
and a reader cannot weigh the positive results without them.

\subsection{Methodological failures we made and corrected}
\label{sec:negative}

\begin{enumerate}
\item \textbf{An invalid ``route is absent'' measurement.} Our first check of
the \texttt{diffpatch} precondition was
\texttt{grep -c diffpatch /usr/local/bin/gitea}, which returned 0. That was
meaningless: the file is a \textbf{581-byte} bash wrapper, and the real binary
is \texttt{/app/gitea/gitea} (124\,MB), where the route strings are present.
Had we not checked file size, we would have published a false negative about
our own feasibility gate.
\item \textbf{A timestamp format that silently truncated.} The canary used
\texttt{date -u +\%Y-\%m-\%dT\%H:\%M:\%S.\%6N\%z}; busybox \texttt{date} does
not expand \texttt{\%N}, so \texttt{executed\_at\_utc} ends as
\texttt{2026-08-27T14:33:22.} (visible in Listing~\ref{lst:canary}). Nothing
is invalidated --- the proof is the planted nonce plus the process lineage,
not the timestamp --- but the field is not usable as sub-second evidence.
\item \textbf{A repository count we do not use.} The impact probe reported
\texttt{total\_repos=0}, which contradicts the same probe's own listing of
two repository paths. The cause is an escaping bug in a hook string of ours
(\texttt{``*.git''} inside a Go-quoted command), not a property of the target.
The count is therefore discarded and no repository total appears anywhere in
this paper.
\item \textbf{An A/B probe we threw away.} A structural probe intended to
record bare/non-bare state and index contents atomically per cycle was
unreliable: the temporary repository lives for milliseconds, our bare check
returned \texttt{error}, and several assertions lost a race against cleanup.
It is \textbf{not} used as evidence. The A/B conclusion rests solely on the
three observations of Table~\ref{tab:ab}.
\item \textbf{A capture that did not exist.} The first \texttt{tcpdump} run
produced no file: without \texttt{sudo} the interface capture fails on
\texttt{CAP\_NET\_RAW}. The wire evidence in this paper was recaptured with
elevated capability (\texttt{evidence/wire-attack.pcap}, 33\,400 bytes, 72
packets). We document it because an unrecorded empty capture is exactly the
kind of evidence that would have been fabricated by omission.
\end{enumerate}

\subsection{Scope limits of the positive results}
\begin{itemize}
\item \textbf{Fix semantics on an official image, not a rebuilt binary.} The
patched arm is the vendor's \texttt{gitea/gitea:1.27.1} image. That is the
strongest available comparison for an operator, but it means our A/B cannot
attribute the outcome to the single hunk of Listing~\ref{lst:fix} alone; 1.27.1
contains other changes. We claim the hunk is decisive from code reading, and
behavioural inertness from observation, not from a single-variable binary
diff.
\item \textbf{The vendor's exfiltration channel is not reproduced.} Our PoC
has no output path by design, so we make no claim about the retrieval branch
technique beyond citing the advisory.
\item \textbf{No availability or privilege-escalation claim.} The write proof
was a single file inside \texttt{/data/gitea}, deleted afterwards; we did not
test container escape, service disruption, or escalation beyond
\texttt{uid=1000(git)}.
\item \textbf{Lab network shape.} Attacker position is the bridge gateway of a
zero-egress lab. Latency, TLS termination, reverse-proxy request buffering and
rate limiting of a production deployment are not measured, and the
$0.82$--$0.89$\,s figure is a lab delivery time, not a real-world one.
\item \textbf{Single lab instantiation per version.} One container per arm,
one configuration per arm. The 22/22 figure is repetition of the attack, not
diversity of deployments.
\end{itemize}

% =============================================================================
\section{Ethics, Reproducibility, and Disclosure Position}
\label{sec:ethics}

\textbf{No discovery claim.} CVE-2026-60004, the mechanism, the vulnerable
command line, and a complete working exploit are the vendor's contribution,
published in GHSA-rcr6-4jqh-j84m on 2026-07-28. This paper states that in the
abstract, in Section~\ref{sec:record}, in the companion video material, and in
the artifact README, so that no reader can mistake a reproduction for a
discovery.

\textbf{Consent and blast radius.} Every experiment ran in a Docker lab we
created and destroyed: \texttt{--internal} bridge with verified zero egress,
fictional hostname \texttt{gitea.lab}, no real users, no real data, no
third-party system. Payloads were canary-first by policy: the maximum side
effect of any proof in this paper is one canary file, one probe file (deleted
after reading), and repository objects inside a container we own. No reverse
shell, no destructive action, no denial of service, no social engineering.

\textbf{Reproducibility.} The companion folder is self-contained:
\texttt{lab/run-lab.sh} brings up both arms and \texttt{lab/verify.sh} asserts
the 16 invariants it checks (internal network flag, both versions reported,
three zero-egress probes --- no route to a public address from either arm and no
name resolution on the vulnerable arm --- plus a host-side routing assertion, the
fictional \texttt{*.lab} hostname, git version floor, executable \texttt{/tmp},
real binary size, presence of the \texttt{diffpatch} string) with a pass/fail
exit code;
\texttt{exploit/cve-2026-60004-canary.py} is a stdlib-only reimplementation;
\texttt{exploit/impact-probe.py} produces the impact table;
\texttt{exploit/ab-monitor.sh} produces the arrival evidence that keeps the
negative arm live; \texttt{evidence/} holds the canary and probe transcripts,
the two PoC run logs, and the wire capture; \texttt{designs/} holds the fact
sheet and the fetched fix diff. Every number in this paper is either in
\texttt{designs/FACTS.md} or in one of those files. A reader with Docker can
reproduce Table~\ref{tab:ab} in one afternoon.

\textbf{Operational note on our own pipeline.} The lab, the payloads, the
measurements and this text were produced by an autonomous agent pipeline under
an operator-approved scope memo, including the negative results of
Section~\ref{sec:negative}. We flag this as a limitation as much as a
convenience: an automated pipeline produced a wrong \texttt{grep} measurement,
a broken timestamp, a bogus repository count and an unusable structural probe,
and each was caught only because the protocol required recording failed
attempts. Pipeline output is trustworthy at exactly the strength of its
negative-result discipline.

% =============================================================================
\section{Conclusion}
\label{sec:conclusion}

CVE-2026-60004 is a composition failure, not a coding error: a patch applied
with \texttt{--index}, a three-way fallback that checks a path out, and a bare
sandbox whose root is \texttt{\$GIT\_DIR}, each defensible alone and lethal
together. The vendor did more than disclose it --- they published the exploit.
What a disclosure cannot supply is the rest of the lifecycle, and that is what
we contribute here: execution proven at process level with the parent named
(\texttt{git write-tree}) and the bare-root equivalence observed
(\texttt{GIT\_DIR=.}), preconditions measured rather than assumed and shown to
hold by default in the official image, a fix verified with a live negative
(22/22 cycles with no execution, arrival positively observed in 2 of 22 samples,
HTTP 201 turning into HTTP 500), impact
quantified instead of listed, and a negative-results log that includes four of
our own broken measurements.

The one thing we would ask a reader to keep is this: \emph{read the fix, not
just the advisory}. Changing \texttt{Clone(\dots,true)} to
\texttt{Clone(\dots,false)} neutralises the exploit without making it
detectable --- the same bytes still land on disk. Any defence built on
recognising the attack's content will fire on benign patches (our mode-bit
signature fired on 3 of 3 legitimate cycles that added an executable file) and
stay silent about whether the host was actually exposed. The defence that works is
behavioural: a \texttt{git} process that spawns a shell.

% =============================================================================
\begin{thebibliography}{99}

\bibitem{ghsa}
Gitea Security Team, ``Remote Code Execution via diffpatch Git Hook
Installation --- GHSA-rcr6-4jqh-j84m (CVE-2026-60004), affected
\texttt{>=1.17, <1.27.1}, patched in 1.27.1, published 2026-07-28, including
a complete Python proof-of-concept,''
\url{https://github.com/advisories/GHSA-rcr6-4jqh-j84m}.

\bibitem{nvd}
NIST NVD, ``CVE-2026-60004'' (status Analyzed; published 2026-08-26;
CVSS 3.1 9.8 Critical, secondary assessment; CWE-94; no CVSS v4.0 score),
\url{https://nvd.nist.gov/vuln/detail/CVE-2026-60004}.

\bibitem{fixcommit}
Gitea, ``refactor: git patch apply (\#38637) (\#38638),'' commit
\texttt{d7bc52be\hspace{0pt}eadff4be5f56\hspace{0pt}90de4d5de4\hspace{0pt}2abd10affe},
2026-07-26, \texttt{services/repository/-files/patch.go}.

\bibitem{kev}
CISA, ``Known Exploited Vulnerabilities Catalog --- CVE-2026-60004, added
2026-08-25, FCEB due date 2026-08-28, catalog version \texttt{2026.08.27},
\texttt{dateAdded} and \texttt{dueDate} reproduced verbatim in the shipped
KEV record, \texttt{knownRansomwareCampaignUse =
Unknown}, notes referencing BOD 26-04,''
\url{https://www.cisa.gov/known-exploited-vulnerabilities-catalog}.

\bibitem{bod}
CISA, ``Binding Operational Directive 26-04'' (staggered remediation windows
of 3, 14 and 60 days referenced by the KEV entry for this CVE).

\bibitem{githooks}
J.~Hamano and the Git project, ``git hooks --- \texttt{post-index-change},''
Git documentation, \url{https://git-scm.com/docs/githooks}.

\bibitem{gitapply}
J.~Hamano and the Git project, ``git-apply --- \texttt{--index},
\texttt{-\/-cached}, \texttt{-3},'' Git documentation,
\url{https://git-scm.com/docs/git-apply}.

\bibitem{cwe94}
MITRE, ``CWE-94: Improper Control of Generation of Code ('Code
Injection'),'' \url{https://cwe.mitre.org/data/definitions/94.html}.

\bibitem{attack}
MITRE, ``MITRE ATT\&CK: T1190, T1059.004, T1083, T1552.001, T1546.004,''
\url{https://attack.mitre.org/}.

\end{thebibliography}

\end{document}
